\UseRawInputEncoding

\lstdefinestyle{promptstyle}{
  basicstyle=\ttfamily\scriptsize,
  breaklines=true,
  breakatwhitespace=false,
  columns=fullflexible,
  keepspaces=true,
  showstringspaces=false,
  tabsize=2
}

\begin{table*}
\begin{tcblisting}{
  listing only,
  breakable,
  enhanced,
  colback=lightblue!5,
  colframe=lightblue!50,
  rounded corners,
  top=6pt,bottom=6pt,left=8pt,right=8pt,boxsep=4pt,
  title=Prompt for Rubric Generation ({\name}),
  listing options={style=promptstyle}
}
Your task is to extract a set of rubric-style instructions from a user's request.
These rubrics will be used as evaluation criteria to check if a response fully meets the request.
Every rubric item must be a universal principle. If any rubric still contains topic-specific references (e.g., names, places, myths, numbers, historical facts), it is automatically invalid.

- **Two Distinct Categories:**
  - [Hard Rule]: Derived strictly from explicit requirements stated in the <request> (format, length, structure, forbidden/required elements, etc.).
  - [Principle]: Derived by abstracting any concrete cues into domain-agnostic quality criteria (e.g., clarity, correctness, sound reasoning, pedagogy).

- **Comprehensiveness:**
  The rubric must cover all critical aspects implied by the request and examples, including explicit requirements and implicit quality standards.

- **Conciseness & Uniqueness:**
  Each rubric must capture a distinct evaluation criterion. Overlapping or redundant criteria must be merged into a single rubric. Wording must be precise and free of repetition.

- **Format Requirements:**
  - Use a numbered list.
  - Each item starts with "The response" phrased in third person.
  - Append [Hard Rule] or [Principle] at the end of each item.
  - Do not include reasoning, explanations, or examples in the final output—only the rubrics.

Here is the request:
{prompt}

Please generate the rubrics for the above request.
\end{tcblisting}
\end{table*}